\section{Lecture 17 (Prashanth Sridhar): Some consequences of the revised King conjecture}

We next explore several consequences of the main theorem of \cite{king} and its proof.



\subsection{Smoothness, properness, and Rouquier dimension of the Cox category}\label{subsec:smooth}

A $\k$-linear triangulated category $\cT$ is called \defi{proper} if $\dim_\k \Hom_{\cT}(T, T') < \infty$ for all $T, T' \in \cT$. There is also a categorical notion of smoothness, but it is probably not helpful to give the definition in a talk: we say $\cT$ is \defi{smooth} if $\cT$ possesses a dg-enhancement whose diagonal bimodule is perfect. Let us not dwell on this latter definition. There are two points to take away here: (1) the category of perfect complexes on a separated scheme $X$ of finite type over $\k$ is smoooth (resp. proper) if and only if $X$ is smooth (resp. proper) (see the proof of \cite{orlov16}*{Proposition 3.31}) and (2) smooth and proper triangulated categories enjoy many of the same properties as derived categories of smooth and proper varieties. One illustration of (2) is the extensive literature on Hodge theory of smooth and proper dg-categories, as originally proposed by Katzarkov-Kontsevich-Pantev~\cite{KKP}.

\begin{thm} Let $X$ be a smooth projective toric variety. The category $D_{\Cox}(X)$ is smooth and proper, and its Rouquier dimension (see Lecture 9 for the definition) is $\dim(X)$.
\end{thm}

\begin{proof}
Choose a smooth toric stack $\X$ as in the definition of $D_{\Cox}(X)$, so that $D_{\Cox}(X)$ is a full triangulated subcategory of $D(\X)$. The stack $\X$ is projective, in the sense of \cite{BLS16}*{Definition 2.5}, and it is tame in the sense of \cite{BLS16}*{Definition 2.1}. It follows from \cite{BLS16}*{Theorem 6.4} that $D(\X)$ is smooth and proper. The properness of $D_{\Cox}(X)$ is an immediate consequence, and smoothness follows from \cite{orlov16}*{Proposition 3.8, Theorem 3.25}. As for Rouquier dimension, recall that there are fully faithful embeddings $D^b(X) \into D_{\Cox}(X) \into D^b(\X)$, and so there are inequalities~$\dim D^b(X) \le \dim D_{\Cox}(X) \le \dim D^b(\X)$. Finally, we have
$$
\dim D^b(X) = \dim{X} = \dim{\X} = \dim D^b(\X),
$$
where the first and third equalities follow from \cite{FH, HHL}.
\end{proof}

\begin{remark}\label{rmk:mythical2}
Loosely speaking, it is a general mantra that smooth and proper triangulated categories should be ``geometric" in some sense. Of course, $D_{\Cox}(X)$ is of geometric origin: it is by definition a subcategory of the derived category of a toric stack. It is an open question as to whether~$D_{\Cox}(X)$, or its invariants, can be given an intrinsic algebro-geometric description.

The intuition is that Cox category is the derived category of a ``mythical" toric variety whose Cox ring is $S$ and whose effective cone and nef cone coincide. Since $D_{\Cox}(X)$ is not closed under tensor products, it is not (monoidally) equivalent to the derived category of a smooth projective scheme or stack, and so it is unlikely that this intuition can be made fully rigorous. But nevertheless, we can ask wheter its~$K$-theory, Hochschild homology, etc. admit geometric interpretations? For example, singularity categories of hypersurfaces do not have a monoidal structure, but their $K$-theory and Hochschild homology nevertheless recover classical geometric invariants of the hypersurface singularity: this has been an active area of research over the past two decades~\cite{BD, BFK1, BFK2, BP, BRTV, BW, BW1,BW2, CT, dyckerhoff, efimov2,  HLP, KP, kim1,pippi, PV,  segal, shk2, shklyarov1}.
\end{remark}

\subsection{Homotopy category description of the Cox category}
\label{subsec:homotopy}
Let $X$ be a smooth projective toric variety with Cox ring $S$, and let $K_{\Theta}(S)$ denote the homotopy category of bounded complexes of $\Cl(X)$-graded free $S$-modules whose terms are direct sums of modules of the form $S(d)$, where $d$ lies in the Bondal-Thompsen collection $\Theta$. By \cite{king}*{Corollary 4.28} (which is used in the proof that $D_{\Cox}(X)$ has a full strong exceptional collection), there is a natural isomorphism
$$
S_{d' - d} \xra{\cong} \Hom_{\Cox}(\OO_{\Cox}(d), \OO_{\Cox}(d'))
$$
for all $d, d' \in \Theta$. It follows that there is a fully faithful functor $\Psi \co K_{\Theta}(S) \to D_{\Cox}(X)$ that sends~$S(d)$ to $\OO_{\Cox}(d)$ for all $d \in \Theta$.

\begin{thm}
The functor $\Psi$ is an equivalence.
\end{thm}

\begin{proof}
Since $D_{\Cox}(X)$ is generated by $\OO_{\Cox}(d)$ for $d \in \Theta$, this is immediate.
\end{proof}


\subsection*{Resolution of the diagonal for the Cox category}
Let $X$ be a smooth projective toric variety with Cox ring $S$.
As emphasized above, a major source of motivation for the revised King's Conjecture is that the HHL resolution of the diagonal of $X$ depends only on the rays of fan determining $X$. That is, the resolution of the diagonal looks the same for every toric stack $\X_i$ corresponding to a maximal chamber of the secondary fan of $X$ (Corollary~\ref{cor:diagonalUniform}). This suggests that the HHL resolution should somehow be an invariant of the Cox category, and the goal of this subsection is to make this intuition precise by constructing a ``resolution of the diagonal for the Cox category" that pushes forward to the HHL resolutions of the stacks $\X_i$.

We begin with an auxiliary construction. Recall that each $X_i$ is of the form $(\Spec(S) \setminus V(B_i)) / G$ for some monomial ideal $B_i \subseteq S$ and an action of a group $G$ arising from the $\Cl(X)$-grading on~$S$. Let $I \ce \sum B_i$, and let $\mathscr{X}$ denote the toric stack $[(\Spec(S) \setminus V(I)) / G]$. We want to take an HHL resolution of the diagonal for $\mathscr{X}$. We cannot quite directly do this, because $\mathscr{X}$ need not be separated (i.e. the diagonal need not be a closed substack), and one can only take an HHL resolution of a closed substack. So instead, we consider the HHL resolution of the closure of the diagonal torus in $\mathscr{X} \times \mathscr{X}$. Let $H$ denote the complex of free $S \otimes_\k S$-modules associated to this complex of $\OO_{\mathscr{X} \times \mathscr{X}}$-modules, and let $\cH_{i, j}$ denote the HHL resolution of the (closure of the) graph of the birational map $\X_i \dashrightarrow \X_j$. The following result is a consequence of technical machinery in~\cite{HHL}; we omit the proof.

\begin{prop}
\label{prop:homotopy}
For any $i, j$, $H|_{\X_i \times \X_j}$
is homotopy equivalent to $\cH_{i, j}$. %Moreover, if both chambers Γi and Γj for Xi and Xj are in the moving
%cone of ΣGKZ(X), then these complexes are equal.
\end{prop}


%Let $\mathcal{R}$ be the HHL resolution of the diagonal for $X$. The complex~$R \ce \bigoplus_{a \in \Cl(X)} \Gamma(X, \mathcal{R})$ is an object in $K_{\Theta}(S \otimes_\k S)$.
Applying the functor $\Psi$ from \S \ref{subsec:homotopy} to $H$ gives an object $\cH\in D_{\Cox}(X \times X)$. Choose a smooth toric stack $\X$ as in the definition of $D_{\Cox}(X)$, so that $D_{\Cox}(X)$ is a full triangulated subcategory of~$D^b(\X)$, and let $\pi_i \co \X \to \X_i$ denote the accompanying morphisms of toric stacks. Of course,~$D_{\Cox}(X \times X)$ is a full triangulated subcategory of $D(\X \times \X)$ as well. It follows from Proposition~\ref{prop:homotopy} that the object~$(\pi_i \times \pi_j)_*\cH$ is homotopy equivalent to $\cH_{i, j}$: see \cite{king}*{Lemma 6.8(1)} for the straightforward proof. Let~$\Phi_{i, j} \co D^b(\X_i) \to D^b(\X_j)$ (resp. $\Phi \co D^b(\X) \to D^b(\X)$ be the Fourier-Mukai transform with kernel~$\cH_{i, j}$ (resp. $\cH$). The functor $\Phi$ induces an endofunctor of $D_{\Cox}(X)$. We have a commutative diagram
\begin{equation}
\label{eqn:phidiagram}
\xymatrix{
D(\X_i) \ar[r]^-{\Phi_{i, j}} \ar[d]_-{\pi_i^*} & D(\X_j) \\
D_{\Cox}(X) \ar[r]^-{\Phi} & D_{\Cox}(X).\ar[u]_-{(\pi_j)_*}
}
\end{equation}
To see this, use that pullbacks and pushforwards may be realized as Fourier-Mukai transforms, and refer to the formula for convolution of Fourier-Mukai kernels~\cite{huybrechts}*{ Proposition 5.10}.
%$(\pi_i)_* \co D_{\Cox}(X) \to D(\X_i)$ denote the functor induced by the (derived) pushforward $D(\X) \to D(\X_i)$ for all the toric stacks $\X_i$ corresponding to maximal chambers in the secondary fan of $X$.



\begin{thm}\label{thm:resDiagCoxVague}
 The endofunctor $\Phi$ of $D_{\Cox}(X)$ is naturally isomorphic to the identity.
\end{thm}

\begin{proof}
One reduces to showing that there is a natural isomorphism $\Phi(\OO_{\Cox}(d)) \cong \Phi(\OO_{\Cox}(d))$ for all $d \in \Theta$, and similarly for morphisms of these objects. This proof is straightforward and can probably be suppressed in this talk: see~\cite{king}*{Proof of Corollary 6.10}.
\end{proof}

\subsection{Noncommutative resolution of singularities}


\begin{defn}\cite{kuznetsov08}
Let $X$ be a scheme. A \defi{categorical (or noncommutative) resolution of singularities for $X$} is a pair of triangulated functors $G \co \cD \to D^b(X)$ and $F \co \Perf(X) \to \cD$, where~$\cD$ is a smooth triangulated category (see \S \ref{subsec:smooth}), such that
\begin{enumerate}
\item We have $\Hom_{D^b(X)}(F(\F), D) \cong \Hom(\F, G(D))$ for all $\F \in \Perf(X)$, $D \in \cD$.
\item There is a natural isomorphism $G F \cong \id_{\Perf(X)}$.
\end{enumerate}
\end{defn}

Noncommutative resolutions of singularities, in the sense of Van den Bergh~\cite{VdB04}, give examples of categorical resolutions of singularities: see \cite{kuznetsov08}*{\S 3}. Affine toric varieties are known to admit noncommutative resolutions of singularities~\cite{spenko-van-den-bergh-inventiones, faber-muller-smith}. The final application of the main result of \cite{king} that we discuss is a generalization of this result beyond the affine case.

Let $X$ be a semiprojective toric variety (see \cite{CLS}*{\S 7.2} for the definition). For instance, any affine or projective toric variety is semiprojective.
Let $G \co D_{\Cox}(X) \to D^b(X)$ be the pushforward functor, and let $F \co \Perf(X) \to D_{\Cox}(X)$ be the pullback. Recall that $F$ is fully faithful, and we established in \S \ref{subsec:smooth} that $D_{\Cox}(X)$ is smooth (the smooth and projective assumptions were not necessary for this argument). Thus, $D_{\Cox}(X)$ is a categorical resolution of singularities for $X$.

So far, we have not applied the main results of \cite{king}, and we do so now. Let $A_{\Theta}$ be the algebra $\End_{D_{\Cox}(X)}(\bigoplus_{d \in \Theta} \OO_{\Cox}(d))$. Since $\bigoplus_{d \in \Theta} \OO_{\Cox}(d)$ is a tilting object for $D_{\Cox}(X)$ (this does not require $X$ to be smooth or projective), we have
$
D_{\Cox}(X) \simeq D^b(A_\Theta).
$
We conclude:

\begin{thm}
\label{thm:ncres}
The category $D^b(A_\Theta)$ is a noncommutative resolution of singularities for $X$. Moreover, $A_\Theta$ has global dimension $\dim{X}$.
\end{thm}

\begin{proof}
We have  proven the first statement, and the second follows from \cite{FH}*{Theorem 2.14}.
\end{proof}

The noncommutative resolutions of singularities for affine toric varieties given by~\cite{spenko-van-den-bergh-inventiones, faber-muller-smith} follow from Theorem~\ref{thm:ncres}, though there is just a bit of additional work needed for this: see \cite{king}*{Corollary 7.13 and Remark 7.14}.
